\section{Rotulación dinámica}

La rotulación dinámica permite superponer gráficos de texto (rótulos de ponente,
títulos de sesión, avisos informativos) sobre el vídeo en directo. El sistema se basa
en imágenes PNG con canal alfa que se inyectan en el pipeline GStreamer mediante el
elemento \texttt{gdkpixbufoverlay}.

\subsection{Implementación}

El módulo \texttt{voctocore/lib/overlay.py} gestiona el ciclo de vida completo de
un rótulo: carga del fichero PNG, fundido de entrada (\textit{blend-in}), mantenimiento
en pantalla y fundido de salida (\textit{blend-out}). El tiempo de fundido se controla
mediante el parámetro \texttt{blend-time} en la sección \texttt{[overlay]} del fichero
de configuración INI, expresado en milisegundos.

La carpeta de rótulos disponibles se define mediante el parámetro \texttt{path} de
la misma sección, lo que permite adaptar la identidad visual a cada evento (logotipos,
tipografías, colores corporativos) sin modificar el código.

\subsection{Integración con la GUI}

La interfaz gráfica expone un selector de fichero PNG y un campo de texto para el
nombre del ponente. El texto se renderiza sobre la plantilla gráfica activa utilizando
la biblioteca Cairo, generando en tiempo real la imagen PNG resultante que se envía
al núcleo mediante el comando \texttt{set\_overlay}.

% --- Figura sugerida ---
% \begin{figure}[H]
%   \centering
%   \includegraphics[width=0.85\textwidth]{figuras/overlay_ejemplo.png}
%   \caption{Fotograma con rótulo de ponente activo sobre la imagen de cámara.}
%   \label{fig:overlay_ejemplo}
% \end{figure}

% --- Figura sugerida ---
% \begin{figure}[H]
%   \centering
%   \includegraphics[width=0.7\textwidth]{figuras/ciclo_vida_overlay.png}
%   \caption{Ciclo de vida de un rótulo: carga, blend-in, visible, blend-out y eliminación.}
%   \label{fig:ciclo_vida_overlay}
% \end{figure}
